\chapter[\emj{🧠}~LRU Cache (Hash Map + Lista Doble)]{\emj{🧠}~LRU Cache (Hash Map + Lista Doblemente Enlazada)}

\begin{dsaquees}

Tu escritorio 🗄️ solo tiene espacio para 5 carpetas. Cada vez que usas
una, la pones ENCIMA (la más reciente). Cuando llega una carpeta nueva y
ya no cabe, tiras la de HASTA ABAJO: la que llevas más tiempo sin tocar
("Least Recently Used", la menos usada recientemente).

Una LRU Cache hace justo eso con datos, y el reto es lograr que TODO
(guardar, buscar, mover a la cima y tirar la vieja) sea instantáneo
$O(1)$. La combinación mágica: un Hash Map para encontrar rápido + una
lista doblemente enlazada para reordenar rápido.

\end{dsaquees}

\begin{dsaotraforma}

Una LRU (Least Recently Used) Cache guarda hasta \verb|capacity| pares
clave-valor y, al llenarse, desaloja el elemento usado hace más tiempo.
Debe hacer \verb|get| y \verb|put| en $O(1)$.

Ninguna estructura sola lo logra, así que se combinan dos:
\begin{itemize}
\item \textbf{Hash Map} \verb|clave -> nodo|: encuentra cualquier
      elemento en $O(1)$.
\item \textbf{Lista doblemente enlazada}: mantiene el orden de uso. El
      frente es el más reciente; el fondo, el candidato a desalojar.
      Mover un nodo o quitarlo es $O(1)$ porque tienes punteros a ambos
      vecinos.
\end{itemize}

Cada acceso (\verb|get| o \verb|put|) mueve el nodo al frente. Al exceder
la capacidad, se elimina el nodo del fondo (y su clave del mapa).

\end{dsaotraforma}

\begin{dsadiagrama}
\begin{verbatim}
Capacidad = 3.  Frente = más reciente,  Fondo = menos reciente.

put(1) put(2) put(3):
  FRENTE [3] <-> [2] <-> [1] FONDO
  mapa: {1:n1, 2:n2, 3:n3}

get(1)  -> mueve 1 al frente:
  FRENTE [1] <-> [3] <-> [2] FONDO

put(4)  -> lleno! desaloja el fondo (2):
  FRENTE [4] <-> [1] <-> [3] FONDO
  mapa: {1,3,4}   (el 2 se eliminó del mapa y de la lista)

Nodos "centinela" (head/tail falsos) evitan casos borde con NULL.
\end{verbatim}
\end{dsadiagrama}

\begin{dsacomofunciona}
\begin{verbatim}
ESTADO tras cada operación (capacidad = 2)   FRENTE=reciente

  put(1,A)   head <-> [1:A] <-> tail          mapa{1}
  put(2,B)   head <-> [2:B] <-> [1:A] <-> tail  mapa{1,2}

  get(1)     → mueve 1 al frente
             head <-> [1:A] <-> [2:B] <-> tail  devuelve A

  put(3,C)   → LLENO: desaloja el del FONDO (2:B, menos reciente)
             head <-> [3:C] <-> [1:A] <-> tail  mapa{1,3}  (2 eliminado)

  get(2)     → no está → devuelve -1

POR QUÉ lista DOBLE:  para borrar el fondo en O(1) necesitas su
'prev'.  El mapa da el nodo en O(1); la lista lo reordena en O(1).

   mapa:  clave → puntero al nodo
          ┌───┐   ┌───┐
   1 ────►│1:A│   │3:C│◄──── 3
          └───┘   └───┘
   (los nodos viven en la lista; el mapa solo apunta a ellos)
\end{verbatim}
\end{dsacomofunciona}

\begin{lstlisting}[language=C++]
CÓMO FUNCIONA (todo O(1))
- Hash Map clave -> puntero al nodo de la lista.
- Lista doble con centinelas head/tail falsos.
- get(k): si existe, mueve su nodo al frente y devuelve el valor.
- put(k,v): si existe, actualiza y mueve al frente; si no, crea al
  frente. Si superas la capacidad, borra el nodo antes del tail.

📝 PSEUDOCÓDIGO
──────────────────────────────────────────────────────────────

get(clave):
    SI clave no en mapa: return -1
    nodo = mapa[clave]; moverAlFrente(nodo)
    return nodo.valor

put(clave, valor):
    SI clave en mapa:
        nodo = mapa[clave]; nodo.valor = valor; moverAlFrente(nodo)
    SINO:
        SI tamaño == capacidad:
            viejo = tail.prev; quitar(viejo); borrar mapa[viejo.clave]
        nodo = nuevo Nodo(clave, valor)
        insertarAlFrente(nodo); mapa[clave] = nodo

💻 CÓDIGO C++
──────────────────────────────────────────────────────────────

#include <unordered_map>
using namespace std;

class LRUCache {
    struct Node { int key, val; Node *prev, *next; };
    int cap;
    unordered_map<int, Node*> mp;
    Node *head, *tail;                 // centinelas

    void remove(Node* n) { n->prev->next = n->next; n->next->prev = n->prev; }
    void addFront(Node* n) {
        n->next = head->next; n->prev = head;
        head->next->prev = n; head->next = n;
    }
public:
    LRUCache(int capacity) : cap(capacity) {
        head = new Node(); tail = new Node();
        head->next = tail; tail->prev = head;
    }
    int get(int key) {
        if (!mp.count(key)) return -1;
        Node* n = mp[key];
        remove(n); addFront(n);        // marcar como reciente
        return n->val;
    }
    void put(int key, int value) {
        if (mp.count(key)) {
            Node* n = mp[key]; n->val = value;
            remove(n); addFront(n);
            return;
        }
        if ((int)mp.size() == cap) {   // desalojar el menos reciente
            Node* lru = tail->prev;
            remove(lru); mp.erase(lru->key); delete lru;
        }
        Node* n = new Node{key, value, nullptr, nullptr};
        addFront(n); mp[key] = n;
    }
};

⏱️ COMPLEJIDAD (Big-O)
──────────────────────────────────────────────────────────────

  Operación   │ Tiempo  │ Espacio
  ────────────┼─────────┼──────────
  get         │ O(1)    │ O(capacity)
  put         │ O(1)    │ O(capacity)

* El O(1) es EL punto del problema; una lista simple daría O(n).
\end{lstlisting}

\begin{dsaentrevistas}

✅ "LRU Cache" (LeetCode 146): pregunta top de Google/Meta/Amazon. La
   respuesta esperada es EXACTAMENTE Hash Map + lista doble. Menos que
   eso no da O(1).

💡 Nodos centinela (dummy head/tail): eliminan todos los casos borde de
   NULL al insertar/borrar. No los omitas: simplifican mucho el código.

⚠️ ¿Por qué lista DOBLE y no simple? Para borrar un nodo en O(1)
   necesitas su nodo anterior. La lista simple obligaría a recorrer O(n).

🔑 Variante LFU (Least Frequently Used): desaloja por frecuencia, no por
   recencia. Más difícil; suele necesitar un mapa extra de frecuencias.
   Menciónala si te la piden.

\end{dsaentrevistas}

\begin{dsaresumen}

• LRU: al llenarse, desaloja el elemento usado hace más tiempo.
• Hash Map (buscar O(1)) + lista doble (reordenar O(1)).
• Frente = más reciente; fondo = candidato a desalojar.
• get/put mueven el nodo al frente; al exceder capacidad, borra el fondo.
• Centinelas head/tail eliminan casos borde con NULL.
• Lista DOBLE es obligatoria para borrar en O(1).
════════════════════════════════════════════════════════════════

\end{dsaresumen}
